\subsection{四高の道型を全尺度で排除する恒等式}
\label{nfh25:section}

\paragraph{確定した更新．}
四rayの整数 row/pair data
$\sum_{j\ne i}p_{ij}=-q_i$，$p_{ij}p_{ji}=q_iq_j-1>0$
では正符号グラフ $P_4$ は存在しない．この排除は正 weights を仮定せず，
$q_i=1$ も用いない．既存の孤立頂点排除と合わせると，三高・四高とも
正 weights の存在は完全マッチング $2K_2$ と同値になる．
さらに weights を仮定しない全符号型も $2K_2$ または $K_3\sqcup K_1$ に限られる．
孤立頂点があれば，整数 gap から gain 一以上の閉路が残り三頂点に生じる．
長さ二は pair product と矛盾するので，残り三頂点は正三角形を作る．
三高で分母一の ray を孤立させることは行和から不可能だから，
三角形が低 ray を含むという従来の結論もこの一つの議論に含まれる．
旧版の「四高の $P_4$ の存在または排除」は解決済みとして研究計画から外す．
この変更だけで大域の三つの残存グラフ型が消えるとは結論しない．

\paragraph{再利用する計算．}
正の道を $1-2-3-4$ とし，正 entries を
$p_{12}=a,p_{21}=A,p_{23}=b,p_{32}=B,p_{34}=c,p_{43}=C$，
負 entries の絶対値を $d,D,e,E,f,F$ と置く．
\[
 q=q_3,\ t=q_1,\ s=bB+1,\ r=cC+1,\ L=B+c,\ D=L+q,
 \quad h=ab-1\ge0,\quad K=BC-1\ge0.
\]
五つの pair と四行から
\begin{gather*}
 T=(a+t)s+qh=aqf,\qquad U=tL+aD+1=De,\\
 V=ts(r+Cq)+q\{aCs+h(r+Cq)+aq\}=qTE
\end{gather*}
を得る．第六 pair は $\mathcal P=UV-DT(tr-q)=0$ を要求するが，
\begin{align*}
 \mathcal P={}&qsK\,t^2
 +\{aqCs(L+D)+q^2hK+aLq^2+s(r+Cq)+Dsq\}t\\
 &+(aD+1)q\{aCs+h(r+Cq)+aq\}+Dq(as+qh)>0.
\end{align*}
展開の鍵は $L(r+Cq)-Dr=qK$ である．
実際の分母を払った残差は
\[
 \mathcal P=q^2aDf\{eE-(q_1q_4-1)\}.
\]
全ての分母は正であり，$s(r+Cq)t>0$ が狭義正値を保証する．
整数性は $ab\ge1$ と $BC\ge1$ に圧縮されるので，この二不等式を満たす
正の実数 data にも排除は適用できる．任意の実数 data への拡張は未主張である．

\paragraph{以後の利用方針．}
三高の低分母一に依存した二つの道型排除は，公開論文では本恒等式に統合できる．
旧証明は必要な出典・独立検算として旧版に残すだけでよい．
新しい局所模型の列挙では，正 weights を要求する四ray非零 core を最初から
三つの標識付き完全マッチングに制限できる．各辺の開区間から有理 weights を
作れるので，この局所段階で実数最適化や閉路探索を繰り返す必要はない．
ただし異なる core で別々に作った weights が大域に一致するとは限らない．
次は残存大域型における matching の相互配置と clutching を調べる．

\paragraph{検証の位置づけ．}
\texttt{verify\_four\_ray\_matching\_v25.py} は Python 標準ライブラリのみで，
元の四行と五 pair による有理式の再構成，最後の pair の残差，
正値恒等式を独立に展開して比較する．旧版の完全マッチング例については
全ての原方程式と全ての basin 不等式を有理数で再検査する．
パラメータの有界 scan は証明に使わない．この機械検算は証明補助であり，
外部査読または proof assistant による形式化ではない．
